\mysection{CRC32計算の例}
\myindex{CRC32}
\label{sec:CRC32}

\newcommand{\URLCRCSRC}{\url{http://go.yurichev.com/17327}}

これは非常に人気のあるテーブルベースのCRC32ハッシュ計算
技法です\footnote{ソースコードはここから取得されました：\URLCRCSRC}。

\lstinputlisting[style=customc]{\CURPATH/CRC.c}

\myindex{\CLanguageElements!for}

私たちは\TT{crc()}関数のみに関心があります。
ちなみに、\TT{for()}文の2つのループ初期化子に注意してください：\TT{hash=len, i=0}。
もちろん、\CCpp 標準はこれを許可しています。
生成されるコードには、1つではなく2つの操作がループ初期化部分に含まれます。

最適化（\Ox）でMSVCでコンパイルしましょう。
簡潔にするため、\TT{crc()}関数のみを私のコメント付きでここにリストします。

\lstinputlisting[style=customasmx86]{\CURPATH/CRC_2_EN.asm}

\Othree オプションでGCC 4.4.1でも同じことを試してみましょう：

\lstinputlisting[style=customasmx86]{\CURPATH/CRC_gcc_O3_EN.asm}

\myindex{x86!\Instructions!NOP}
\myindex{x86!\Instructions!LEA}

GCCは\NOP と \TT{lea esi, [esi+0]}を追加することで、
ループの開始を8バイト境界に整列しました
（これも\IT{アイドル操作}です）。
これについての詳細は、npadセクション（\myref{sec:npad}）で読むことができます。